% =========================================================
%  Section : Le théorème de Rokhlin
% =========================================================

\subsection{Rokhlin(1952)}

Le théorème de Rokhlin, démontré en 1952, est la première obstruction différentielle connue
à la réalisation d'une forme d'intersection en dimension $4$. Il concerne les variétés spin,
c'est-à-dire, comme rappelé précédemment, les variétés dont la seconde classe de
Stiefel-Whitney s'annule~: $w_2(TM) = 0$. Cette condition est équivalente, pour une variété
simplement connexe, à ce que la forme d'intersection soit paire.

\begin{theorem}[Rokhlin, 1952]\label{thm:rokhlin}
  Soit $M$ une $4$-variété lisse compacte orientée fermée spin. Alors
  \[
    \sigma(M) \equiv 0 \pmod{16}.
\]
\end{theorem}

La contrainte algébrique habituelle — que la signature d'une forme paire unimodulaire est
divisible par $8$ — est ainsi doublée dès que la variété admet une structure lisse. C'est
précisément cet écart entre $8$ et $16$ qui est d'origine différentielle.

\subsubsection*{Applications et limites}

La conséquence immédiate pour les formes d'intersection est la suivante~:

\begin{corollary}\label{cor:rokhlin_e8}
  Il n'existe pas de $4$-variété lisse compacte orientée fermée spin dont la forme
  d'intersection est $Q_{E_8}$.
\end{corollary}

\begin{proof}
  On a $\sigma(E_8) = 8 \not\equiv 0 \pmod{16}$.
\end{proof}

Plus généralement, toute forme paire unimodulaire dont la signature est congrue à $8$ modulo
$16$ est exclue. Cela élimine par exemple $E_8$, $E_8 \oplus E_8 \oplus E_8$ (signature
$24$), et plus généralement $E_8^{\oplus (2k+1)}$ pour tout $k \geq 0$.

\medskip

Le théorème de Rokhlin a cependant une limite importante~: il ne s'applique qu'aux variétés
spin. Si $w_2(TM) \neq 0$, la forme d'intersection de $M$ est impaire, et le théorème ne
donne aucune information. En particulier, il ne dit rien sur les formes définies positives
impaires comme $\langle 1 \rangle^{\oplus n}$, ni sur les variétés non spin à forme paire mais
indéfinie.

\medskip

C'est précisément pour lever cette restriction que le théorème de Donaldson, présenté dans la
section suivante, représente un saut qualitatif majeur~: il fournit une obstruction à la
lissabilité valable pour toute forme définie positive, sans aucune hypothèse sur la structure
spin, en faisant intervenir des méthodes radicalement différentes issues de la géométrie des
équations de Yang-Mills.